\section{The dependency frontier}\label{sec:frontier}

This section states the exact mathematical state of the proof in this paper. Every result of Sections~\ref{sec:criterion} to~\ref{sec:proofA} is in one of the following classes.
\begin{enumerate}
\item[(P)] Proved here, with a complete argument in the paper.
\item[(X)] Reduced exactly to a published theorem, cited where it is used.
\item[(O)] Internal and open. The statement is a Hypothesis of the paper, or a step that a proof marks as not given. The formal development proves it, and this paper does not yet give a conventional proof.
\item[(U)] A cited published theorem whose statement has not been compared with its source, or whose source differs from the form used.
\item[(F)] A statement that fails.
\end{enumerate}
No statement of the paper is in class (F). The reconstruction found no statement that is false, no circular dependency and no hypothesis that contradicts another.

We name each open statement by its number in this paper, and we also give the identifier of the corresponding statement in the formal development, so that a reader can locate it there. The identifiers serve only to locate statements, and the formal development plays no role in the arguments of the paper.

\subsection{The conditional theorem}

Sections~\ref{sec:criterion} to~\ref{sec:reduction} prove Theorem~\ref{thm:intro} from Theorem~A, with the published analytic inputs (C1) to (C9), Linnik's theorem and the bounds of Rosser and Schoenfeld. Section~\ref{sec:proofA} proves Theorem~A from two statements, Hypothesis~\ref{r6:R6:thm:selmer} (the Selmer bound on windows, which gives condition (D6)) and Hypothesis~\ref{r6:R6:thm:records} (the existence of chain records). Hypothesis~\ref{r6:R6:thm:records} collects the results of Sections~\ref{sec:lifts} to~\ref{sec:twisted}, and each of those results is proved from the hypotheses of its section. The open statements on which Theorem~\ref{thm:intro} rests are therefore the hypotheses listed in Subsections~\ref{fr:sec:chain} to~\ref{fr:sec:assembly}, together with the published theorems of Subsection~\ref{fr:sec:audit}.

The following results of Part~II are in class (P), apart from the published theorems that they cite.
\begin{enumerate}[leftmargin=*]
\item Section~\ref{sec:chain}. The covariance of the theta kernels.
\item Section~\ref{sec:inert}. The theta function of a CM curve at an inert prime and the kernel law up to the constant $c$.
\item Section~\ref{sec:gain}. The classification of semilinear pairings and the valuation estimates $g^A$ and $g^B$.
\item Section~\ref{sec:chaindiv}. The divisibility of the composite of a chain datum and the commutation of Hecke operators with the chain.
\item Section~\ref{sec:lifts}. The structure of the lifts, the bound (D1) for denominators, the identification of the kernels with their finite Fourier transforms, the Schr\"odinger model at a split orbit, the integrality of the Hecke commutation coefficients, the passage from entries to lattices at the even positions, and the orbit bounds given (N1) to (N3).
\item Section~\ref{r2:sec:loss}. The reduction of Hecke nondegeneracy to one edge, the colength theorem, Theorems~\ref{r2:A2} and~\ref{r2:A1}, the window theorem, and the analytic and arithmetic tools of the three engines.
\item Section~\ref{r3:sec}. The frame theorem, and Theorem~\ref{r3:A3} from its inputs (a) to (f).
\item Section~\ref{sec:bessel}. The unfolding of Bessel Poincar\'e series, the bad weight at one prime, the pair bound, and Theorems~\ref{w2:S}, \ref{w2:I} and~\ref{w2:B} from their hypotheses.
\item Section~\ref{r4:sec:R4}. The mean period formula, the height bound, the deductions of \eqref{A4}, \eqref{A6} and \eqref{A7} from their hypotheses, the kernel norm theorem, the branching bound and the base-side bound.
\item Section~\ref{r5:sec:localterms}. The bound for the local terms, and \eqref{A5} from (E1) to (E6).
\item Section~\ref{sec:twisted}. The comparison of Tate--Shafarevich groups with class groups, the reduction of the Selmer bound to two statements, the reduction of (V) to (OneChar) and (Count), and steps one and two of the congruence bound.
\item Section~\ref{sec:proofA}. The construction of the nodes, the deduction of \eqref{A1} to \eqref{A7} from a chain record, and the bound $e_1+\dots+e_5+Ae_6+e_7\le x^{3\varepsilon}\mathcal B$.
\end{enumerate}

\subsection{Corrections made in this version}

The reconstruction changed the following statements of earlier versions. Theorem~A is stated with a base scale and a Selmer condition, for windows $[Y,Y^{3/2}]$ with $Y$ arbitrarily large, and its constant is $C=3$. The constant $\log\Omega$ is $\gamma+\log(2\pi/\mathrm{rd}_E)-\lambda_1$, which differs from the period $\gamma+\log(2\pi/\sqrt d)-\lambda_1$ of the CM curve by $\log(\mathrm{rd}_E/\sqrt d)$. The Hecke operators of the chain have the normalizing exponents $n^j_{v,i}=N'_v(j,i)$, and Section~\ref{sec:chaindiv} applies the one-lift relation to the raw operators. Slopes and heights are twice those of Bost in Sections~\ref{r3:sec} and~\ref{r4:sec:R4}, which use the same definitions. The local terms $E_u$ of \eqref{A5} are integrality defects of the lifts at the auxiliary and level primes and do not involve Galois cohomology, while the congruence defects at the twisted places enter \eqref{A3} and \eqref{A4}. The bottom weights $b_t$ of the nodes are chosen together with the weights $a_t$. The covolume bound uses the arithmetic Hilbert--Samuel theorem for each model of a tower of models and an exact identity at the good primes, without a uniform version of that theorem.

\subsection{The chain of lifts}\label{fr:sec:chain}

\begin{enumerate}[leftmargin=*]
\item \emph{(N1), Hypothesis~\ref{r1:hyp:H1}.} The $p$-adic valuation of the frame factor of the step kernel \eqref{stepkernel} at an orbit is $\mu_v(j+1)$ times the change of lengths, up to units. It feeds Proposition~\ref{r1:prop:S3}, hence the one-lift relation, Theorem~\ref{r1:thm:Cii} and Theorem~\ref{r1:thm:R1main}. It needs a precise definition of the frame factor through the Hecke character $\mu$ of Proposition~\ref{covariance}(iii) and a computation from its infinity type. Formal development, the last clause of \texttt{AwayTie}. Class (O), a normalization computation.
\item \emph{(N2), Hypothesis~\ref{r1:hyp:H2}.} The passage from pairings with products of minors to the integral structures at the odd positions. It feeds Corollary~\ref{r1:cor:ent} and the orbit bounds. Formal development, \texttt{adm7}, through \texttt{pending\_tuple\_admissible\_of\_shape}. Class (O). The mechanism is the transfer at the even positions (Theorem~\ref{r1:thm:ent-even}) applied to the contragredient model, with Jacobi's theorem on complementary minors.
\item \emph{(N3), Hypothesis~\ref{r1:hyp:H3}.} The constant $c$ of Theorem~\ref{kernellaw} for the normalized kernel series satisfies $\lvert c\rvert\le1$, which fixes the normalization of the base vector $\theta_0$. It feeds the orbit bounds, and it is also an input of Hypothesis~\ref{w3:W3:K3}. Formal development, \texttt{adm2b}, from \texttt{K1\_representatives\_LT\_of\_dictionary}. Class (O), a normalization statement.
\end{enumerate}
The published inputs of Section~\ref{sec:lifts} are the bideterminant basis of Weyl modules (Green, De Concini--Eisenbud--Procesi), Minkowski's bound for finite subgroups of $\mathrm{GL}_n(\mathbb Z)$, Borel's finiteness of class numbers and Jacobowitz's classification of unimodular hermitian lattices. They are in class (U), because the comparison with their sources is not complete and the audit did not locate the statement of Jacobowitz used for unramified places.

\subsection{Hecke nondegeneracy and the loss}\label{fr:sec:loss}

\begin{enumerate}[leftmargin=*]
\item \emph{Hypothesis~\ref{r2:engines}, the engines of Hecke nondegeneracy.} Formal development, the three engines behind \texttt{LPprime}. Through Theorem~\ref{r2:LPthm} and Corollary~\ref{r2:meankappa} it feeds the colength bound and so \eqref{A1} and \eqref{A2}. Subsections~\ref{w1:sec:C}, \ref{w1:sec:B} and~\ref{w1:sec:A} prove the window theorem, the tie identity and the window bound at one level, the factorial digit split, the decay of the unit parts, the divided-derivative congruence, the existence and commutation of Fock projectors, the bound for the extremal projector on dominant weight spaces, the lifting from the residue field, the Frobenius identity, the digit decomposition, and the log-binomial and frame-deficit estimates. The remaining open steps are the statements (C6) to (C11) of engine~(C), (B7) to (B9) of engine~(B), and (A2), (A5), (A6), (A7), (A9) of engine~(A). Of these, (C8), (C9), (C11), (B7), (B8), (B9), (A5) and (A7) need several pages each, and the others are short computations. The outline of the proof of Theorem~\ref{r2:LPthm} also leaves the Pieri splitting at the tops to be written out. Class (O). No statement here has a published counterpart, and none is in doubt. The norm bound (B7) for the isotypic projector is the step most likely to change a constant, and a loss of larger order than $r^2\log_pa$ per projector would have to be checked against the size conditions (Z).
\item \emph{Inputs (L1) to (L9) and the size conditions (Z) of Subsection~\ref{r2:sec:lossinputs}.} They are definitions and size bounds given by Sections~\ref{sec:lifts}, \ref{r5:sec:localterms} and~\ref{sec:proofA}. The comparison exponent $e_p$ of (L2), the colength at $T_u$ of (L7) and the bound $\Pi\le\exp(x^{\varepsilon/4})$ in (Z1) are not written out. Class (O), interface.
\end{enumerate}
Published inputs: the integral standard monomial bases of Doubilet--Rota--Stein and De Concini--Eisenbud--Procesi, the extremal projector of Asherova--Smirnov--Tolstoy, the surjectivity of multiplication of Ramanan--Ramanathan, the theorems of Kummer and Legendre, Krull's intersection theorem and the elementary divisor theorem. The first three are in class (U).

\subsection{The slope inequality and Bessel Poincar\'e series}\label{fr:sec:frame}

Theorem~\ref{r3:A3} proves \eqref{A3} from inputs (a) to (f). Input (a) needs the divisibility of the chain and the definition of $N_u$, which Sections~\ref{sec:lifts} and~\ref{r2:sec:loss} give. Input (b) is Corollary~\ref{r4:R4:htcor}, input (c) is the bound $s+\mu_{\max}(R)\le x^{C_H}b$ of the construction, and inputs (d) and (e) are Theorems~\ref{w2:B}, \ref{w2:I} and~\ref{w2:S}. Section~\ref{sec:bessel} proves these three theorems from the following statements. Formal development, \texttt{BesselId}, \texttt{LinkIRR} and \texttt{LinkS}.
\begin{enumerate}[leftmargin=*]
\item \emph{Hypothesis~\ref{w2:H-see}, the seesaw relation.} Every linear relation among the periods of the chain on the definite side transfers to the local Bessel periods of the theta lifts on the indefinite side. Its proof needs the chain of seesaw identities of the theta tower, the anisotropic Siegel--Weil formula, Siegel Eisenstein series and the local doubling zeta integrals that give the constants. Class (O), a major statement. It is the theta-tower seesaw relation behind the Bessel identity.
\item \emph{Hypotheses~\ref{w2:H-dict} and~\ref{w2:H-rat}.} Dictionaries between the abstract Bessel realization and the automorphic objects, and the rationality of the period maps. Class (O), interface statements of the construction, with the rapid decay of cusp forms as a published input.
\item \emph{Hypotheses~\ref{w2:H-cnt} and~\ref{w2:H-DIM}, the multiplicity formula.} The number of eigenforms with labels in a given set equals $D^*$ times a Plancherel weight up to a relative error $\exp(-cNf\log N)$, and $D_{\mathrm{all}}$ is close to $D^*$. Its mechanism is a pre-trace formula for a projector with a discrete-series matrix coefficient, with the identity term giving $D^*$ and the other terms bounded by reduction theory, filtrations of Grayson--Stuhler type and volume ratios of Siegel--Weil type. Class (O), a major statement. It feeds both Theorem~\ref{w2:I} and Theorem~\ref{w2:S}.
\item \emph{Hypothesis~\ref{w2:H-KE}, the kernel expansion.} The absolute convergence of the unfolding of the pairings of Poincar\'e series and its organization into the term family of the pair bound. Class (O), several pages.
\item \emph{Hypothesis~\ref{w2:H-MR}, the main row.} The main-row estimate, a ratio of masses through the Siegel--Weil formula in its convergent range and Schur orthogonality. Class (O), a major statement.
\item \emph{Hypothesis~\ref{w2:H-LR}, the loop rows.} The loop rows are exponentially small, from the congruence gain at the typed primes against polynomial bounds for the archimedean amplitudes, summed over coranks with genus means of lattice counts. Class (O), a major statement.
\item \emph{Hypotheses~\ref{w2:H-NG} and~\ref{w2:H-DG}.} The nongeneric row and the degenerate terms. Class (O), one estimate each.
\item \emph{Hypotheses~\ref{w2:H-wt} and~\ref{w2:H-rows}.} The comparison of the level idempotents with Gelfand--Graev idempotents over a finite field, and the identification of the rows of $R$ with the good forms. Class (O), a finite group computation and an interface statement.
\item \emph{Hypothesis~\ref{w2:H-SZ}, the sizes, and the window \eqref{w2:window}.} Class (O), interface with Section~\ref{sec:proofA}.
\end{enumerate}
The kernel expansion, the multiplicity formula and the main-row and loop-row estimates are the statements of the paper with the largest amount of unwritten mathematics. Their mechanisms are classical in type (Petersson-type unfolding, pre-trace formula, Siegel--Weil, Schur orthogonality, lattice counting), and their explicit dependence on the parameters is particular to this proof. The archimedean inputs, uniqueness of Bessel functionals, integrability of the Bessel functions of discrete series, Schur orthogonality with the formal degree and polynomial bounds for amplitudes, are of published type (Harish-Chandra, Schmid, Wallach, Jiang--Sun--Zhu), with an explicit dependence on the gap of the Harish-Chandra parameter that the paper does not write out.

\subsection{Heights and the target slope}\label{fr:sec:heights}

Theorems~\ref{r4:R4:A4}, \ref{r4:R4:A6} and~\ref{r4:R4:A7} prove \eqref{A4}, \eqref{A6} and \eqref{A7} from Hypotheses~\ref{r4:R4:HK}, \ref{r4:R4:HP}, \ref{r4:R4:HL}, \ref{r4:R4:HT}, \ref{r4:R4:HW} and~\ref{r4:R4:HD}. Subsections~\ref{w3:sec:HK} to~\ref{w3:sec:HT} reduce the first four further.
\begin{enumerate}[leftmargin=*]
\item \emph{Kernel norms.} Theorem~\ref{w3:W3:HK} proves Hypothesis~\ref{r4:R4:HK} from Hypotheses~\ref{w3:W3:K3} (the normalized kernel, which depends on (N3)), \ref{w3:W3:K4} (the mean of the logarithm of the theta-null over conjugates) and~\ref{w3:W3:K5} (class and index counts). Class (O), three short statements, the second resting on the integrality and the conjugate bounds of Weber function values, which are in class (U).
\item \emph{Final pairing.} Theorem~\ref{w3:W3:HP} proves Hypothesis~\ref{r4:R4:HP} from Hypothesis~\ref{w3:W3:xi}, an inequality between Fock norms of products of minors, and from the factorization of the final pairing through the branching map, which the paper does not write out. Class (O), short.
\item \emph{Frame constants.} Proposition~\ref{w3:W3:HLprop} reduces Hypothesis~\ref{r4:R4:HL} to the Rallis inner product formula in the convergent range and to Harish-Chandra's formula for the formal degree, for the theta lifts of the chain, together with the identification of the maps $\theta_j$ with the period maps of Theorem~\ref{w2:B}. The two formulas are published in type and in class (U). Their application to the chain is class (O).
\item \emph{Target slope.} The assembly in Subsection~\ref{w3:sec:HT} proves Hypothesis~\ref{r4:R4:HT} from the base-side bound (proved), the passage inequality (proved from Hypothesis~\ref{w3:W3:mult}, two multiplicity counts), the covolume bound Hypothesis~\ref{w3:W3:L3}, and Theorems~\ref{r5:thm:Ebound} and~\ref{r5:thm:DIVinN} for the local terms. Hypothesis~\ref{w3:W3:L3} rests on Hypothesis~\ref{w3:W3:unif}, the uniform mean slope on deep covers, whose argument uses a tower of integral models indexed by a prime cutoff, the arithmetic Hilbert--Samuel theorem on each model, Zhang's theorem on successive minima and an exact identity at the good primes. Class (O). The arithmetic Hilbert--Samuel theorem for a fixed model and Zhang's lower bound are in class (X). The form of Zhang's theorem used drops the ampleness of the generic fibre, and a perturbation argument restores it. A uniform arithmetic Hilbert--Samuel theorem for the family of covers would also suffice, and no such theorem is published. The exactness of the identity at the good primes is the point on which the uniformity in the level rests.
\item \emph{Hypotheses~\ref{r4:R4:HW} and~\ref{r4:R4:HD}.} The weight relations of the kernels and the lower bound for the dimension budget $D^*$, from Siegel's mass formula for unimodular hermitian lattices and Harish-Chandra's formal degree. Class (O) for the first, (U) for the published inputs of the second.
\end{enumerate}
The mean period formula, Proposition~\ref{r4:R4:period}, is in class (P) from Kronecker's first limit formula, Jacobi's product formula, Dirichlet's class number formula, the main theorem of complex multiplication and the discriminant formula for a tower of fields. The formula of Chowla and Selberg in the $\sqrt d$ form agrees with it after the correction by $\log(\mathrm{rd}_E/\sqrt d)$.

\subsection{The local terms}\label{fr:sec:local}

Corollary~\ref{r5:cor:A5} proves \eqref{A5} with $g=\log(g_0g_1)$ and $e_5=C_EmB_0$ from Hypothesis~\ref{r5:prop:H1} (the degree identity of the lifts and size bounds) and Hypothesis~\ref{r5:prop:H6I} (the inert orbit losses). Both are computations with the definitions of Sections~\ref{sec:lifts}, \ref{r2:sec:loss} and~\ref{sec:proofA}. Class (O), small. No Galois cohomology enters \eqref{A5}.

\subsection{The twisted places}\label{fr:sec:twisted}

\begin{enumerate}[leftmargin=*]
\item \emph{Hypothesis~\ref{w4:thm:vwin}, the Selmer bound on windows.} Formal development, \texttt{VWin}. Theorem~\ref{w4:thm:decomp} reduces it to (I), the plus part with the prime $2$ and the bad primes on windows, and (II), the minus part. The generic plus part reduces to (CorGK), which reduces through (PropGK) to (Count) and (OneChar) (Propositions~\ref{w4:prop:corToBound} to~\ref{w4:prop:onechar2prop}). (OneChar) rests on (G1) to (G4) of Subsection~\ref{tw:sec:vandiver}. Class (O).
\begin{enumerate}
\item The unresolved core is the automorphy assertion (G3) inside (OneChar), the finiteness of the ordinary deformation ring with residual representation $\Psi\oplus1\oplus\Psi^{-1}$, and automorphy of every ordinary point that is irreducible on open subgroups. It is not a published theorem, because hypotheses (4), (5) and (6c) of Theorem~1.1 of Allen--Newton--Thorne~\cite{ant20} fail in this configuration. Formal development, \texttt{GdEndAlt}, inside \texttt{OneChar}. The paper does not claim a proof of (G3).
\item The consequence (V) of Subsection~\ref{tw:sec:vandiver}, and the bound of Proposition~\ref{w4:prop:need} that Hypothesis~\ref{w4:thm:vwin} implies by itself, are stronger than the published bounds in the direction of Vandiver's conjecture of Soul\'e~\cite{soule99} and de~la~Torre~Dur\'an~\cite{delatorre}, by one exponential in the threshold. The formal development proves (V) as \texttt{SmallClassNumbers.Bound}. This is the one place where an internal statement of the proof appears to contain mathematics stronger than the available published literature.
\item (G2), an intersection inequality on a graded Rees algebra, is commutative algebra with no arithmetic input. (G4) has the shape of Theorem~A of Newton--Thorne~\cite{nt23} at irreducible points and of Kato's theorem~\cite{kato04} at reducible points, both in class (U) for the form used. (G1) is Ribet's construction. (Count) is a polynomial count of algebraic automorphic forms on a definite unitary group in three variables with purity.
\item The part of (I) at $l=2$ and at the bad primes uses torsion bounds of the type of Soul\'e and Dwyer--Friedlander, and the minus part (II) is of the type of the Main Conjecture for abelian fields with the theorem of Ferrero--Washington. Both are class (O) as written proofs, with published inputs in class (U).
\item The normalization of the exponent $e(x)$ differs between Hypothesis~\ref{w4:thm:vwin} and Hypothesis~\ref{r6:R6:thm:selmer} by the integer $c$ of Proposition~\ref{field}(iii). Class (O), interface.
\end{enumerate}
\item \emph{Hypothesis~\ref{w4:thm:linkn}, the congruence bound.} Formal development, \texttt{LinkN}. Steps one and two are proved (Propositions~\ref{w4:prop:cong} and~\ref{w4:prop:fitting}). The open parts are the identification of $\eta_j(w)$ with the module of Proposition~\ref{w4:prop:cong}, the descent along a map of Hecke algebras, conditions (F2) and (F3) for the blocks that occur, and step three, which consists of the block bounds, the lifting theorem for crystalline classes with Wach modules at the places above $l$, the twist, the rationality and integrality of twisted theta series, and the final sums. Class (O). The lifting theorem with Wach modules and the statements on twisted theta series are the largest parts.
\item \emph{The equality $d^0_j=D_0$} used in Theorem~\ref{r5:thm:DIVinN}. Class (O), small.
\end{enumerate}

\subsection{The assembly}\label{fr:sec:assembly}

\begin{enumerate}[leftmargin=*]
\item \emph{Hypothesis~\ref{r6:R6:thm:records}, chain records.} The inequalities (K1) to (K8) are stated with the exponents $e_D=2$ and $e_L=0$ of the formal development. Sections~\ref{sec:lifts} to~\ref{sec:twisted} prove inequalities of the same shape from their hypotheses. The check that their constants and exponents give (K1) to (K8) with absolute exponents, and the construction of the words, the auxiliary primes $T_u$ with their congruences, the lattices $M_u$ and $R$ and the period maps, is not written out. Class (O), interface. Formal development, \texttt{linkCore} with \texttt{LinkG}, \texttt{LinkS}, \texttt{LinkIRR}, \texttt{LinkDIV}, \texttt{LinkMu}, \texttt{LinkN}.
\item \emph{Hypothesis~\ref{r6:R6:thm:selmer}}, the form of Hypothesis~\ref{w4:thm:vwin} that gives (D6).
\end{enumerate}

\subsection{Published theorems}\label{fr:sec:audit}

The comparison of the cited published theorems with their sources has the following state.
\begin{enumerate}[leftmargin=*]
\item \emph{Verified against the source.} Berger's theorem that de Rham representations are potentially semistable~\cite{berger02}, Th\'eor\`eme~0.7 (Corollaire~5.22), whose proof uses the theorem of Andr\'e, Kedlaya and Mebkhout. The formula of Chowla and Selberg in the $\sqrt d$ form. The arithmetic Hilbert--Samuel theorem for a fixed model, and Zhang's lower bound, with the perturbation argument above. The theorems of Barnet-Lamb, Gee, Geraghty and Taylor on Galois representations for regular algebraic polarized cuspidal representations over CM fields~\cite{blggt14} and the local-global compatibility of Caraiani, in the normalizations of~\cite{blggt14}, with the theorem of Gelfand and Kazhdan on contragredients and the tensor product theorem of Flath as inputs whose sources were not located. The geometric lemma and the structure of the Grothendieck ring of Bernstein and Zelevinsky. Henniart's characterization of the local Langlands correspondence. M\'inguez's theorem on Howe duality for type II pairs.
\item \emph{Repairable differences.} The derivation of (C7) from Gallagher's theorem needs an orthogonality step, which the comparison wrote out, and the printed statement of Gallagher's Theorem~7 still has to be transcribed. A theorem on irreducibility of products of unitary representations is stated for all non-archimedean local fields and was verified in characteristic zero, which is the case used. The tensor product theorem of Flath is used for unitary representations without an admissibility hypothesis. The local-global statement for Galois representations should be cited from~\cite{blggt14} and Caraiani's paper, and not from Harris--Taylor alone.
\item \emph{Not located in this pass.} The uniform prime ideal theorem (C9) for ray classes of $\mathbb Q(\sqrt{-d})$ in the Landau--Page range. Results of Moeglin--Waldspurger and Harish-Chandra on cusp forms (rapid decay, finiteness, discrete series). The Langlands classification of irreducible representations of $\mathrm{GL}_n$ in the quotient form. Clozel's results on purity and cohomological representations. The normalization of local invariants in Poitou--Tate duality. The theory of Bloch and Kato, Wach modules, the Herr complex, Mazur's presentation lemma, Taylor--Wiles--Kisin patching, Thorne's theorem for nice primes, Kato's theorem in the range used, adjoint vanishing of Newton--Thorne type, Kolyvagin systems of Gauss sums, Soul\'e's torsion exponent. Siegel--Weil and doubling in the forms used, the Rallis inner product formula in the range $n<m\le2n$, Howe duality for real groups, Siegel's mass formula for hermitian lattices, Jacobowitz's classification. Weber function values. Explicit constants in the theorem of Angehrn and Siu, which have to be derived from its proof.
\item \emph{A sign to check.} The sign convention for the Hodge--Tate weights of the $l$-adic characters attached to algebraic Hecke characters, which enters many uses of these characters.
\end{enumerate}
No published theorem was found to be misquoted in a way that breaks a step.

\subsection{The order of the next steps}

The statements above group into the following problems, in the order in which they bear on the validity of the argument.
\begin{enumerate}
\item The automorphy statement (G3), with (G2), (G4) and (Count), and the remaining parts (I) and (II) of the Selmer bound. This is the only place where the proof appears to go beyond the published literature, and it needs specialists in Hida theory, residually reducible patching and Iwasawa theory.
\item The theta-tower seesaw relation, the multiplicity formula, the kernel expansion and the main-row and loop-row estimates of Section~\ref{sec:bessel}.
\item The congruence bound, with the lifting theorem for crystalline classes and the statements on twisted theta series.
\item The uniform mean slope on deep covers and the covolume bound, with the exactness of the identity at the good primes.
\item The engines of Hecke nondegeneracy.
\item The normalization statements (N1) to (N3), the short statements of Section~\ref{r4:sec:R4}, the local terms, and the interface statements of Section~\ref{sec:proofA}.
\item The comparison of the remaining published theorems with their sources, starting with (C9), the results on cusp forms and Siegel--Weil, and the Hodge--Tate sign.
\end{enumerate}
